% subsection: 階比型 \, $a_{n+1}=f(n)a_n$

  \subsection{\texorpdfstring{階比型 \, $a_{n+1}=f(n)a_n$}{階比型}}\label{節：階比型}
  \begin{bookpoint}
    \textbf{倍率を掛ける}\par
    添字ごとの倍率を掛け合わせる方法と,その倍率に合う因子で割る方法を比べよう.零で割らない条件は共通である.
  \end{bookpoint}
    \bookterm{difference}{階差}型では変化を足した.ここでは,各段階の倍率を掛け合わせる.
    難しい点は,積を計算することと,倍率をそろえる因子を見つけることに分かれる.
    まず積の表示を学び,その後に\bookterm{normalize}{正規化}の因子を選ぶ理由を見る.

    具体的な説明の前に1つだけ記号を新たに導入しておく.
    適当な数が並べられた\bookterm{sequence}{数列}$x_n$が存在した時,
    \[
      \prod_{i=1}^n x_i
      =x_1\times x_2\times\cdots\times x_n
    \]
    となるような記号$\prod$（ラージパイ）を定義する.
    これを\bookterm{product}{総積}記号といい,積を表す.
    また,一般に以下が成り立つ.
    \begin{align*}
      &\prod_{i=1}^{n}i=n!\\
      &\prod_{i=1}^{n}k=k^n\\
      &\prod_{i=1}^{n}a_ib_i=\displaystyle\prod_{i=1}^na_i\displaystyle\prod_{i=1}^nb_i\\
      &\prod_{i=1}^{n}a_i=\displaystyle\prod_{i=1+k}^{n+k}a_{i-k}\\
      &\prod_{i=1}^{n}(i+k)=\dfrac{(n+k)!}{k!}
    \end{align*}
    ここでの$a_i$や$b_i$は\bookterm{recurrence}{漸化式}等によって定義されるような,法則性を持った\bookterm{sequence}{数列}である必要はなく,
    「1から順に番号が振られた,数の並び」ぐらいのイメージで問題ない.
    では再び, \,$a_n$の\bookterm{general}{一般項}を求めていく.
    \begin{align*}
      &a_{n+1}=f(n)a_n\\
      &a_{n+1}=f(n)\{f(n-1)a_{n-1}\}\\
      &a_{n+1}=f(n)\left[f(n-1)\{f(n-2)a_{n-2}\}\right]\\
      &\qquad \vdots\\
      &a_{n+1}=f(1)f(2) \cdots f(n)a_1\\
      &a_{n+1}=a_1 \prod_{k=1}^n f(k)\\
      &\therefore \quad a_n=a_1 \prod_{k=1}^{n-1} f(k)
    \end{align*}
    この積による表示は$f(k)=0$となる場合も成り立つ(空積は$1$とする).
    \noindent \bookblocklink{example-ratio-1}{problem}{answer}{【例題】}\par\nobreak
    \begin{enumerate}[label=(\arabic*),parsep=3pt,format=\bookitemlink{example-ratio-1}{problem}{answer}]
      \item $\displaystyle a_1=1,\, a_{n+1}=2na_n$
      \item $\displaystyle a_1=2,\, a_{n+1}=\frac{n+2}{n}a_n$
      \item $\displaystyle a_1=1,\, a_{n+1}=n(n+1)a_n$
    \end{enumerate}

    \noindent\bookblocklink{example-ratio-1}{hint}{problem}{【方針】}\par\nobreak
    \begin{enumerate}[label=(\arabic*),parsep=3pt,format=\bookitemlink{example-ratio-1}{hint}{problem}]
      \item 係数$2n$のうち$n$を階乗で吸収すれば,一定の倍率$2$だけが残ります.
      \item $n(n+1)$を一つのまとまりとして,隣の添字にずらした比を計算してみましょう.倍率$\frac{n+2}{n}$が現れます.
      \item ちょっと頭を捻ってみましょう.先程までは$n!$で割りましたが,今回はうまくいくでしょうか.\,$f(n)=n(n+1)$の形に合わせて,どのように階乗を組み合わせれば「きれい」になるか考えてみてください.
    \end{enumerate}

    \noindent\bookblocklink{example-ratio-1}{answer}{problem}{【解答】}\par\nobreak
    \begin{enumerate}[label=(\arabic*),parsep=3pt,format=\bookitemlink{example-ratio-1}{answer}{problem}]
      \item $\displaystyle a_n=(n-1)!\cdot 2^{n-1}$
      \item $\displaystyle a_n=n(n+1)$
      \item $\displaystyle a_n=n!(n-1)!$
    \end{enumerate}

    \noindent\bookblocklink{example-ratio-1}{solution}{problem}{【解法】}\par\nobreak
    \begin{enumerate}[label=(\arabic*),parsep=3pt,format=\bookitemlink{example-ratio-1}{solution}{problem}]
      \item $\displaystyle a_1=1,\, a_{n+1}=2na_n$\\
            両辺を$n!$で割って,
            \begin{align*}
              &a_{n+1}=2na_n\\
              &\frac{a_{n+1}}{n!}=2\frac{a_n}{(n-1)!}
            \end{align*}
            ここで
            \[
              b_n=\frac{a_n}{(n-1)!}
            \]
            と置くと,
            \begin{align*}
              &b_{n+1}=2b_n\\
              &b_n=b_1\cdot 2^{n-1}
            \end{align*}
            $b_n$を$a_n$に戻して
            \begin{align*}
              &\frac{a_n}{(n-1)!}=\frac{a_1}{(1-1)!}2^{n-1}\\
              &\therefore \quad a_n=a_1 \cdot (n-1)! \cdot 2^{n-1}
            \end{align*}

      \item $\displaystyle a_1=2,\, a_{n+1}=\frac{n+2}{n}a_n$
            \begin{align*}
              &a_{n+1}=\frac{n+2}{n}a_n\\
              &na_{n+1}=(n+2)a_n\\
            \end{align*}
            両辺を
            \[
              n(n+1)(n+2)
            \]
            で割って,
            \begin{align*}
              &\frac{a_{n+1}}{(n+1)(n+2)}=\frac{a_{n}}{n(n+1)}
            \end{align*}
            $a_n$について遡っていくと
            \begin{align*}
              &\frac{a_{n}}{n(n+1)}=\frac{a_{n-1}}{(n-1)(n)}=\\
              &\frac{a_{n-2}}{(n-2)(n-1)}=\cdots=\frac{a_{1}}{1(1+1)}\\
              &\therefore \quad \frac{a_{n}}{n(n+1)}=\frac{a_{1}}{1(1+1)}\\
              &a_n=n(n+1)
            \end{align*}
      \item $\displaystyle a_1=1,\, a_{n+1}=n(n+1)a_n$\\
              両辺を$n!(n+1)!$で割って,
            \begin{align*}
              &a_{n+1}=n(n+1)a_n\\
              &\frac{a_{n+1}}{n!(n+1)!}=\frac{n(n+1)a_n}{n!(n+1)!}\\
              &\frac{a_{n+1}}{n!(n+1)!}=\frac{a_n}{n!(n-1)!}\\
              &\therefore \quad \frac{a_n}{n!(n-1)!}=\frac{a_1}{1!(1-1)!}\\
              &a_n=n!(n-1)!
            \end{align*}
    \end{enumerate}

    因子を見つけにくければ,先に倍率の積を書き,隣り合う分子と分母が相殺するかを調べるとよい.
    \bookterm{normalize}{正規化}はその積の計算を,一歩ずつの更新として表したものである.

    なお,ここでは\bookterm{ratio}{階比}型\bookterm{recurrence}{漸化式}$a_{n+1}=f(n)a_n$について,\,$f(n)$が指数関数でない場合について扱ったが,
    より複雑な関数であったり指数関数を用いたりする場合は\bookterm{logarithm}{対数}を用いることがある.
    それら次節応用7パターンで解説する.


\bookstep{掛ける量を選び、倍率をそろえる}\label{解釈：階比と正規化}
    一方,以下では非零の因子を掛けて倍率をそろえる.これは\sectionref*{節：補助数列・正規化}の\bookterm{normalize}{正規化}に当たる.
    $f(k)\ne0$を仮定し,常に$g(n)\ne0$となるように選ぶ.
    探すのは,\,$\frac{g(n+1)}{g(n)}$が元の倍率$f(n)$の逆数になる因子である.
    与式の両辺に$g(n+1)$を掛けると,
    \[
      g(n+1)a_{n+1}=g(n+1)f(n)a_n
    \]
    ここで$g(n+1)f(n)=g(n)$を満たすような$g(n)$を選べば,
    \begin{align*}
      &g(n+1)a_{n+1}=g(n)a_n\\
      &g(n+1)a_{n+1}=g(n)a_n=g(n-1)a_{n-1}\\
      &=\cdots=g(1)a_1\\
      &\therefore \quad a_n=\frac{g(1)a_1}{g(n)}
    \end{align*}
    この解法では,\,$b_n=g(n)a_n$の更新が$b_{n+1}=b_n$となる.
    項ごとに変わる倍率を取り除いたため,初項の値だけを追えばよくなったのである.
    これを背景に考えると,例えば$f(n)=n$のときは$g(n)=\dfrac1{(n-1)!}$と取れば$g(n+1)f(n)=g(n)$が成り立ち,
    \begin{align*}
      &a_{n+1}=na_n\\
      &\frac{a_{n+1}}{n!}=\frac{a_n}{(n-1)!}\\
      &\frac{a_n}{(n-1)!}=\frac{a_1}{0!}\\
      &\therefore \quad a_n=(n-1)!\cdot a_1
    \end{align*}
    が言える.
    多くはこの場合が出題されるため,演習問題はこの手法を使って解くものを揃えた.
